\section{\texorpdfstring{\(J\)}{J}-normalization at discrepancy one}
\label{sec:normalization}

Let \(\phi\colon Y\to Z\) be the flag bundle of \(V\oplus V'\), so that
\(\phi^{*}V\) and \(\phi^{*}V'\) carry filtrations whose successive quotients
are line bundles, and \(\phi^{*}\colon H^{*}(Z)\to H^{*}(Y)\) is injective;
this is the splitting principle for extremal quantum cohomology,
\cite[Corollary~4.3]{SS}.  The Chern roots \(\rho_{i}\) and \(\sigma_{j}\)
exist individually on \(Y\), and we continue to write \(V\), \(V'\),
\(\rho_{i}\), \(\sigma_{j}\) for the pullbacks.  Write
\[
  A=H^{*}(\PP_{Y}(V))
   \cong H^{*}(Y)[H]\big/\left\langle
     \textstyle\prod_{i=1}^{r}(\rho_i+H)\right\rangle
\]
for the coefficient ring of \eqref{eq:ss35} after this pullback.  It is a
finite-dimensional graded \(\C\)-algebra whose elements of positive degree are
nilpotent.  Each factor of the denominator of \eqref{eq:ss35} has the
form \(mz+a\) with \(m\ge1\) and \(a=\rho_{i}+H\in A\), so
\[
  \frac1{mz+a}
  =\frac1{mz}\sum_{l\ge0}\Bigl(\frac{-a}{mz}\Bigr)^{l}\in z^{-1}A[z^{-1}]
\]
is a finite sum, so each summand of \eqref{eq:ss35} lies in
\(zA[z^{-1}][[q]]\) once the \(\ln q\) carried by the prefactor is handled as
described before Lemma~\ref{lem:jslice}.
By the \emph{\(z\)-order} of such an element we mean the largest \(n\) with a
nonzero \(z^{n}\) coefficient.  A \(z\)-order bound proved after the pullback
is a coefficientwise statement about classes in the image of an injective
map, so it descends to the original cohomology ring; the bound below is in
any case stated in terms of \(r\), \(s\) and \(d\) alone.

\begin{lemma}
  \label{lem:order}
  For every \(d\ge1\), the degree-\(d\) summand of \eqref{eq:ss35} has
  \(z\)-order at most
  \[
    1+s(d-1)-rd=1-s-\nu d .
  \]
\end{lemma}

\begin{proof}
  The numerator \(\prod_{j=1}^{s}\prod_{m=0}^{d-1}(\sigma_{j}-H-mz)\) has
  \(sd\) factors.  The \(s\) factors with \(m=0\) are \(\sigma_{j}-H\) and
  contain no \(z\); each of the remaining \(s(d-1)\) factors is linear in
  \(z\).  So the numerator is a polynomial in \(z\) over \(A\) of degree at
  most \(s(d-1)\).  The denominator \(\prod_{i=1}^{r}\prod_{m=1}^{d}(\rho_{i}
  +H+mz)\) has \(rd\) factors, all with \(m\ge1\), so by the expansion above
  its inverse lies in \(z^{-rd}A[z^{-1}]\).  Multiplying by the leading
  factor \(z\) of \eqref{eq:ss35} adds one.
\end{proof}

The two hypotheses of Theorem~\ref{thm:normalization} enter here and nowhere
else.

\begin{corollary}
  \label{cor:normalized}
  If \(\nu=1\) and \(s\ge1\), then every summand of \eqref{eq:ss35} with
  \(d\ge1\) has \(z\)-order at most \(1-s-d\le-1\), and hence
  \[
    z\,e^{t/z}q^{c_{1}(T)/z}
    \sum_{d\ge0}q^{(r-s)d}\,
    \frac{\prod_{j=1}^{s}\prod_{m=0}^{d-1}(\sigma_{j}-H-mz)}
         {\prod_{i=1}^{r}\prod_{m=1}^{d}(\rho_{i}+H+mz)}
    =z\one+\bar{\mathbf t}+O(z^{-1}),
  \]
  with \(\bar{\mathbf t}=t+\ln(q)c_{1}(T)\) as in \eqref{eq:barparam}.
\end{corollary}

\begin{proof}
  The bound of Lemma~\ref{lem:order} is \(1-s-d\), which is at most \(-1\)
  for \(d\ge1\) and \(s\ge1\).  The \(d=0\) summand is
  \(z\,e^{t/z}q^{c_{1}(T)/z}=z\,e^{\bar{\mathbf t}/z}
  =z\one+\bar{\mathbf t}+O(z^{-1})\).
\end{proof}

\begin{remark}
  \label{rem:boundary}
  Both hypotheses are sharp.  For \(\nu=1\) and \(s=0\), so \((r,s)=(1,0)\),
  the bound of Lemma~\ref{lem:order} is \(1-d\) and is attained by the formal
  \(d=1\) term of the hypergeometric expression:
  \[
    z\cdot\frac{q}{\rho_{1}+H+z}
    =q\one-\frac{q(\rho_{1}+H)}{z}+\cdots.
  \]
  Its \(z^{0}\) coefficient is \(q\one\), so this formal expression is not
  \(J\)-normalized.  Geometrically \(\PP(V)=Z\) has point fibres and no line
  class \([L]\), so \((1,0)\) is a degenerate endpoint of the extremal-line
  setup, not a further standard-flip case repaired here.  For \(\nu=0\) the bound is the
  \(d\)-independent value \(1-s\) and the sum is not \(q\)-adically
  convergent, consistently with the standing assumption \(r>s\) in \cite{SS}.
\end{remark}

To convert the normalization into an identification we use the standard
characterization of the \(J\)-function as a slice of Givental's cone.  Let
\(\mathcal L_{T}\) denote the overruled Lagrangian cone of the genus-zero
Gromov--Witten theory of \(T\), in the extremal specialization
\eqref{eq:barparam}.

The parameter \(\bar{\mathbf t}=t+\ln(q)c_{1}(T)\) is not literally defined
over \(\C[[q]]\).  Two readings remove the difficulty, and they agree.  One
may adjoin \(\ln q\) to the coefficients, the standard extension used with
the divisor equation, over which both the extremal cone and \eqref{eq:ss35}
are defined; or one may factor off the prefactor \(q^{c_{1}(T)/z}\) and work
with the remaining series, in which \(q\) appears only in integral powers.
In either reading the degree-\(d\) summand carries \(q^{(r-s)d}\) by
\eqref{eq:ss35}.  That weight is positive because \(r>s\), so the series
converges \(q\)-adically and \eqref{eq:ss35} is a point of the ambient
symplectic space.

\begin{lemma}[\(J\)-slice uniqueness; {\cite{Givental}}, {\cite{CoatesGivental}}]
  \label{lem:jslice}
  Let \(f\in\mathcal L_{T}\) be of the form \(f=z\one+\tau+O(z^{-1})\) for
  some \(\tau\) in the parameter space.  Then \(f=J^{T}(\tau,z)\).
\end{lemma}

This is the mechanism behind every mirror theorem of Givental type: a point
of the cone with a prescribed \(z^{0}\) coefficient is the \(J\)-function at
that coefficient, and the mirror map is precisely the difference between
the prescribed \(z^{0}\) coefficient and the parameter one started with.

The second ingredient is that \eqref{eq:ss35} lies on \(\mathcal L_{T}\) at
all.  This is asserted in \cite[Remark~4.5(3)]{SS}, the same remark
whose other assertion Corollary~\ref{cor:normalized} contradicts, so we prove
it instead of quoting it.  We begin with split bundles, where the two sources
the authors name do the work.

\begin{lemma}
  \label{lem:cone}
  Let \(r>s\ge0\) and suppose \(V=\bigoplus_{i=1}^{r}L_{i}\) and
  \(V'=\bigoplus_{j=1}^{s}M_{j}\).  Then the right-hand side of
  \eqref{eq:ss35} lies on \(\mathcal L_{T}\).
\end{lemma}

\begin{proof}
  The projective bundle \(\PP(V)\to Z\) is a toric fibration with fibre
  \(\PP^{r-1}\).  Brown's theorem \cite[Theorem~1]{Brown} states that the
  hypergeometric modification of the \(J\)-function of the base lies on the
  overruled cone of the total space, and \cite[Corollary~1]{Brown} is the
  projective-fibration case.  Setting the Novikov variables of \(Z\) to zero
  and keeping the fibrewise variable is the extremal specialization
  \cite[(33)]{SS}, a homomorphism of Novikov rings by
  \cite[Proposition~4.1]{SS}; it leaves \cite[(34)]{SS}, which therefore
  lies on the extremal cone \(\mathcal L_{\PP(V)}\) of \(\PP(V)\).  Brown's
  hypotheses concern the fibration
  \(\PP(V)\to Z\) and the Chern roots \(\rho_{1},\dots,\rho_{r}\); no second
  bundle occurs in them, so they cannot constrain \(r-s\).

  Now the twist.  A stable map \(f\colon C\to\PP(V)\) of class \(k[L]\) with
  \(k\ge1\) has image in a fibre of \(\PP(V)\to Z\) by \cite[Lemma~2.6]{SS}.
  There \(V'(-1)=\pi^{*}V'\otimes\mathcal O_{\PP(V)}(-1)\) pulls back to a
  sum of \(s\) line bundles of degree \(-k_{i}\) on each component
  \(C_{i}\), where \(k_{i}\ge0\) is the degree of \(f|_{C_{i}}\) and
  \(\sum_{i}k_{i}=k\).  A section vanishes on every component with
  \(k_{i}>0\), and, the dual graph being a connected tree with at least one
  such component, on the contracted components as well; so
  \(H^{0}(C,f^{*}V'(-1))=0\) and
  \(V'(-1)\) is concave along every extremal class, while its dual
  \(E:=(V'(-1))^{\vee}=\pi^{*}(V')^{\vee}\otimes\mathcal O_{\PP(V)}(1)\) is
  convex there.  Twisting by the inverse equivariant Euler class of a concave
  bundle gives the Gromov--Witten theory of its total space, and the
  hypergeometric modification \cite[Theorem~2]{CoatesGivental} of the
  \(J\)-function has a Serre-dual partner obtained by replacing \(e,E\) with
  \(e^{-1},E^{\vee}\); both are the discussion ``On Serre duality'' of
  \cite{CoatesGivental}.

  The dual modification is termwise, so its degree-\(d\) factor can be read
  off.  The Chern roots of \(E\) are \(H-\sigma_{j}\), and each pairs to
  \(d\) on \(d[L]\), so the factor contributed by the \(j\)-th root is
  \(\prod_{k=-d+1}^{0}\bigl(-\lambda-(H-\sigma_{j})+kz\bigr)\), which is
  empty for \(d=0\).  In positive degree the concave local theory admits the
  non-equivariant specialization, and letting \(\lambda\to0\) and writing
  \(m=-k\) turns this into
  \(\prod_{m=0}^{d-1}(\sigma_{j}-H-mz)\).  Applying the modification to
  Brown's projective-bundle \(I\)-function \cite[(34)]{SS}, and passing to
  the extremal specialization appropriate to \(T\), which replaces
  \(q^{rd}\) by \(q^{(r-s)d}\) and \(c_{1}(\PP(V))\) by \(c_{1}(T)\), gives
  \eqref{eq:ss35}.  The \(d=0\) summand is untouched by the modification and
  remains the classical term \(z\,e^{\bar{\mathbf t}/z}\).  The modification
  involves the Chern roots of \(V'(-1)\) and nothing else; here too \(r\) and
  \(s\) enter only as numbers of factors.
\end{proof}

The passage from split bundles to general ones is where the printed
attribution of \cite[Remark~4.5(3)]{SS} cannot be followed at \(\nu=1\), so
we spell it out.  What it transports is the identity between
\eqref{eq:ss35} and \(J^{T}\), not membership in an infinite-dimensional
cone.

\begin{lemma}
  \label{lem:split-reduction}
  Let \(r>s\ge0\).  If the right-hand side of \eqref{eq:ss35} equals the
  extremal \(J\)-function for every local model whose bundles split, then it
  does so for every local model.
\end{lemma}

\begin{proof}
  Let \(\phi\colon Y\to Z\) be the flag bundle of \(V\oplus V'\) and
  \(\tilde\phi\colon T_{Y}=Y\times_{Z}T\to T\) the induced map.  The
  Cartesian square identifies the relative obstruction theories and hence
  the virtual classes of \(\overline M_{0,n}(T_{Y},d)\) and
  \(\overline M_{0,n}(T,d)\) for \(d\) in the span of \([L]\); this is the
  content of \cite[Theorem~4.2]{SS} and, for the fundamental solution, of
  \cite[Lemma~9.5]{SS}.  Applied to the two-point descendant invariants of
  \cite[(18)]{SS} it gives \(\tilde\phi^{*}S^{T}=S^{T_{Y}}\tilde\phi^{*}\)
  and therefore \(\tilde\phi^{*}J^{T}=J^{T_{Y}}\).  Both sides of
  \eqref{eq:ss35} are symmetric in the Chern roots
  \cite[Remark~4.5(2)]{SS}, and the expression for \(T_{Y}\) is the pullback
  of the expression for \(T\).  Since \(\tilde\phi^{*}\) is injective by
  Leray--Hirsch \cite[Corollary~4.3]{SS}, it suffices to prove the identity
  on \(T_{Y}\), where \(V\) and \(V'\) are filtered with line-bundle
  quotients.

  It remains to replace a filtered bundle by its associated graded.  Take one
  step \(0\to A\to V\to Q\to0\) of the filtration.  The construction of
  \cite[Lemma~9.4]{SS} produces a bundle \(\ker(g)\) on \(Z\times\mathbf
  A^{1}\) with \(\ker(g)_{1}=V\) and \(\ker(g)_{0}=A\oplus Q\), hence a
  smooth projective family \(\PP(\ker g)\to Z\times\mathbf A^{1}\) carrying
  the bundle \(V'(-1)\).  By \cite[Proposition~3.1]{SS} the moduli spaces
  \(\overline M_{0,n}(\PP(\ker g)_{t},k[L])\) form a smooth proper family
  over \(\mathbf A^{1}\), and concavity makes
  \(R^{1}\pi_{\mathcal C*}f^{*}V'(-1)\) a relative locally free obstruction
  bundle of rank \(s(k-1)\).  Its Euler class defines a relative virtual
  class restricting on each fibre to the extremal virtual class of
  \cite[Corollary~3.2]{SS}.  Deformation invariance therefore identifies all
  extremal genus-zero descendant invariants, and with them the extremal
  \(J\)-functions, under the isomorphism of cohomology rings of
  \cite[Lemma~9.4]{SS}.  The right-hand side of \eqref{eq:ss35} is likewise
  unchanged, because \(V\) and \(A\oplus Q\) have the same Chern roots.
  Repeating the construction for the filtration of \(V'\), and iterating both
  filtrations as in \cite[Lemma~9.6]{SS}, reaches the split
  associated-graded local model.
\end{proof}

\begin{proposition}
  \label{prop:IJ}
  Let \(r=s+1\) with \(s\ge1\).  Then the right-hand side of
  \eqref{eq:ss35} equals the extremal \(J\)-function \(J^{T}(q,t,z)\) of the
  local model, in the variables and completion of \cite[Theorem~4.4]{SS}.
\end{proposition}

\begin{proof}
  Assume first that \(V\) and \(V'\) split.  By Lemma~\ref{lem:cone} the
  right-hand side of \eqref{eq:ss35} is a point of \(\mathcal L_{T}\), and by
  Corollary~\ref{cor:normalized} its expansion is
  \(z\one+\bar{\mathbf t}+O(z^{-1})\).  So it is a point of \(\mathcal
  L_{T}\) with prescribed \(z^{0}\) coefficient \(\bar{\mathbf t}\), and by
  uniqueness of the \(J\)-slice, Lemma~\ref{lem:jslice} with
  \(\tau=\bar{\mathbf t}\), it is \(J^{T}(q,t,z)\).  The general case follows
  from Lemma~\ref{lem:split-reduction}.
\end{proof}

\begin{remark}
  \label{rem:circularity}
  The route just taken differs from the printed attribution of
  \cite[Remark~4.5(3)]{SS} in one place.  The authors deduce cone membership
  from \cite{Brown}, \cite{CoatesGivental} and \cite[Lemma~9.6]{SS}.  The
  first two are the split-case inputs of Lemma~\ref{lem:cone} and carry no
  hypothesis on \(\nu\).  But \cite[Lemma~9.6]{SS} rests on
  \cite[Lemma~9.4]{SS}, whose identification of the fundamental solutions of
  \(T\) and its associated-graded model is deduced from \eqref{eq:ss35}
  read as a \(J\)-function formula, that is, from \cite[Theorem~4.4]{SS} with
  its hypothesis \(r-s>1\).  At \(\nu=1\) that formula is what is being
  established, so Lemma~\ref{lem:split-reduction} appeals to the underlying
  deformation instead of to the conclusion drawn from it.  Once
  Proposition~\ref{prop:IJ} is available, \cite[Lemmas~9.4--9.6]{SS} hold at
  \(\nu=1\) as printed, their input having been supplied; this is what
  Section~\ref{sec:scope} uses.
\end{remark}

\begin{corollary}
  \label{cor:presentation}
  Let \(r=s+1\) with \(s\ge1\).  Then the presentation \eqref{eq:ss37} of
  \cite[Theorem~4.6]{SS} holds, the modified extremal \(J\)-function
  \cite[(39)--(40)]{SS} of \(T\) is as printed, and
  \cite[Theorem~1.2]{SS} holds: after the extremal
  specialization, \(c_{1}(X)\star_{q}-\) on \(QH^{*}(X)\) has eigenvalue
  \(0\) with multiplicity \(\rk H^{*}(X')\) and the single nonzero
  eigenvalue
  \[
    \lambda_{0}(q)=e^{-\pi\sqrt{-1}\,s}q=(-1)^{s}q
  \]
  with multiplicity \(\rk H^{*}(Z)\).
\end{corollary}

\begin{proof}
  The proof of \cite[Theorem~4.6]{SS} specializes the parameter \(t\) of
  \cite[Theorem~4.4]{SS}, exhibits a differential operator annihilating
  \(J^{T}\), and concludes with \cite[Theorem~10.3.1]{CoxKatz} and
  \cite[Lemma~2]{Pandharipande}; it uses \(r-s>1\) only to know that
  \eqref{eq:ss35} is \(J^{T}\).  Proposition~\ref{prop:IJ} supplies that,
  so the proof applies verbatim, and \cite[(39)--(40)]{SS} follow by the
  same substitution as in \cite[Section~4.3]{SS}.

  For the spectrum, \cite[Propositions~5.1 and~5.2]{SS} concern the abstract
  ring \eqref{eq:ss37} and its subspaces and assume only \(r>s\), and
  \cite[Corollary~5.3]{SS} combines them with the ring homomorphism
  \(i^{*}\colon QH^{*}_{\mathrm{ext}}(X)\to QH^{*}_{\mathrm{ext}}(T)\) of
  \cite[Theorem~3.3]{SS}.  Once \eqref{eq:ss37} is available at \(\nu=1\),
  that argument is unchanged.  Its output is \(r-s\) nonzero eigenvalues
  \(\lambda_{k}(q)=(r-s)e^{-\pi\sqrt{-1}(2k+s)/(r-s)}q\) for
  \(0\le k<r-s\); at \(r-s=1\) only \(k=0\) occurs, and
  \(\lambda_{0}(q)=e^{-\pi\sqrt{-1}s}q\).
\end{proof}

\begin{remark}
  \label{rem:pbundle}
  The companion statement \cite[(34)]{SS} for \(\PP(V)\) is the \(s=0\) case
  of the same expression, and Lemma~\ref{lem:order} applied to it gives
  maximal \(z\)-order \(1-rd\).  This is at most \(-1\) for every \(d\ge1\)
  as soon as \(r\ge2\), which holds whenever \(r=s+1\) with \(s\ge1\).  So
  \cite[(34), (36) and~(38)]{SS} are covered in the same range, by the same
  count.
\end{remark}

Proposition~\ref{prop:IJ} and Corollary~\ref{cor:presentation} prove
Theorem~\ref{thm:normalization}.  They also settle the status of
\cite[Remark~4.5(3)]{SS}: under the standing assumption \(r>s\) its range
\(r-s\le1\) is exactly \(\nu=1\), and by
Corollary~\ref{cor:normalized} and Remark~\ref{rem:boundary} its negative
assertion holds there only for \(s=0\).
